\chapter{Claude Code: A Practitioner's Reference}
\label{app:claude-code}

\begin{remark}[Purpose of this Appendix]
This appendix is a comprehensive reference for Claude Code, the agentic
command-line interface developed by Anthropic that was used to draft, score,
refine, and maintain this book.  Readers who wish to understand the tool's
architecture, configure it for their own research workflows, or build on the
skills and agents defined in \texttt{.claude/} will find the material here
essential.  The appendix follows the natural onboarding arc: what the tool is,
how to install and configure it, the built-in tools and extension mechanisms,
multi-agent orchestration, and finally patterns that are particularly productive
for quantitative finance research.
\end{remark}

% ---------------------------------------------------------------
\section{What is Claude Code}
\label{app:claude-code:overview}

\begin{context}
On 27 February 2025, Anthropic released a research preview of a terminal-native
coding assistant named Claude Code.  The reception was unusual: engineers who
tried it reported that it was the first AI coding tool that could take a vague
description of a task and actually complete it across dozens of files without
constant hand-holding.  Within thirteen months, 4\% of all public GitHub
commits---roughly 135{,}000 per day---were attributed to Claude Code, a
42{,}896$\times$ growth rate.  By early 2026, Anthropic reported that
90\% of its own production code was AI-written, predominantly via Claude Code.
The tool had graduated from curiosity to infrastructure.
\end{context}

\subsection{Overview and Positioning}
\label{app:claude-code:positioning}

Claude Code is an \emph{agentic} coding assistant that runs inside a terminal
(or an IDE terminal panel) and has persistent access to the developer's file
system, shell, and external services.  Unlike chat-based assistants that operate
purely in a conversation window, Claude Code can read files, write files, run
shell commands, spawn background processes, call external APIs via
Model-Context-Protocol servers, and coordinate fleets of parallel sub-agents.

The key architectural decision that distinguishes Claude Code from earlier AI
coding tools is that the model is placed \emph{in the loop of the build process}
rather than acting as a side-channel advisor.  When the model writes a function,
it can immediately compile it, run the test suite, read the error output, and
iterate---all within a single conversation turn if the user's permission settings
allow it.

\subsection{Supported Models}
\label{app:claude-code:models}

Claude Code works with the full range of production Anthropic models.  As of
mid-2026, the three actively supported model tiers are:

\begin{itemize}
  \item \textbf{Claude Opus 4.8} --- the flagship reasoning model; required for
        Agent Teams (Section~\ref{app:claude-code:agent-teams}) and
        recommended for complex multi-file refactors.
  \item \textbf{Claude Sonnet 4.6} --- the default model; balances capability
        and cost for day-to-day coding tasks and lecture drafting.
  \item \textbf{Claude Haiku 4.5} --- the fast, low-cost model; useful for
        high-frequency sub-agent calls (e.g., batch file classification) where
        latency and cost dominate accuracy concerns.
\end{itemize}

The active model can be overridden per session with the \texttt{/model} command
or set permanently in \texttt{settings.json}.  Individual sub-agents and
Workflow stages can specify their own model via the \texttt{model} option,
enabling cost-aware tiering within a single orchestration run.

\subsection{Key Milestones}
\label{app:claude-code:milestones}

The evolution of Claude Code's feature set matters for practitioners because
community resources and blog posts are often written against older versions.
Table~\ref{tab:claude-code-milestones} summarises the major capability
releases.

\begin{table}[ht]
\centering
\caption{Major Claude Code capability releases (2024--2026).}
\label{tab:claude-code-milestones}
\begin{tabular}{lll}
\toprule
\textbf{Date} & \textbf{Feature} & \textbf{Notes} \\
\midrule
November 2024 & MCP support & Model Context Protocol; first external tool bridge \\
February 2025 & Research preview & Public launch with Bash, Read, Edit, Write tools \\
July 2025     & Subagents (Agent tool) & Spawning isolated child sessions \\
September 2025 & Hooks & Pre/post tool-call shell scripts \\
October 2025  & Skills \& Plugins & Markdown slash commands; bundled plugin packages \\
October 2025  & Workflow tool & Deterministic multi-agent fan-out scripts \\
February 2026 & Agent Teams & Peer-to-peer multi-session coordination \\
March 2026    & Scheduled Routines & Cron-style remote execution on Anthropic infra \\
\bottomrule
\end{tabular}
\end{table}

% ---------------------------------------------------------------
\section{Installation and Setup}
\label{app:claude-code:setup}

\subsection{Prerequisites and Authentication}
\label{app:claude-code:install}

Claude Code requires Node.js 18 or later.  Installation is a single
\texttt{npm} command:

\begin{minted}{bash}
npm install -g @anthropic-ai/claude-code
\end{minted}

On first launch, Claude Code opens a browser-based OAuth flow to authenticate
against an Anthropic account.  For headless servers (CI/CD environments, remote
compute nodes), set the \texttt{ANTHROPIC\_API\_KEY} environment variable
instead; the browser flow is then skipped.

\begin{minted}{bash}
export ANTHROPIC_API_KEY="sk-ant-..."
claude                        # start interactive session
claude "fix all type errors"  # non-interactive single-shot mode
\end{minted}

To run Claude Code with a specific model from the command line:

\begin{minted}{bash}
claude --model claude-opus-4-8
\end{minted}

\subsection{IDE Integration}
\label{app:claude-code:ide}

Claude Code ships first-party IDE extensions for VS Code and JetBrains IDEs
(IntelliJ IDEA, PyCharm, WebStorm, etc.).  The extensions embed the full
Claude Code terminal panel inside the IDE, with additional GUI affordances:

\begin{itemize}
  \item \textbf{Inline diffs} --- proposed file edits are shown as unified diffs
        inside the editor gutter before the user approves them.
  \item \textbf{Plan review pane} --- in Plan mode
        (Section~\ref{app:claude-code:plan-mode}), the plan document is
        rendered in a dedicated editor panel and can be edited directly.
  \item \textbf{\texttt{@}-mention} --- any open file can be dragged into the
        conversation with \texttt{@filename} to provide immediate context.
  \item \textbf{Conversation history} --- previous sessions are browsable via
        the IDE's activity bar.
\end{itemize}

Install the VS Code extension from the Marketplace by searching
\emph{Claude Code} (publisher: \texttt{Anthropic}).  For the integrated
terminal to support multi-line input via \textsc{Shift+Enter}, run the
following inside a Claude Code session once after installation:

\begin{minted}{text}
/terminal-setup
\end{minted}

This writes the required keybinding to VS Code's \texttt{keybindings.json}.

% ---------------------------------------------------------------
\section{Core Interface}
\label{app:claude-code:interface}

\subsection{CLI Commands and Navigation}
\label{app:claude-code:cli}

Claude Code's slash commands are the primary navigation layer.
Table~\ref{tab:claude-code-slash} lists the most commonly used built-in
commands.

\begin{table}[ht]
\centering
\caption{Built-in Claude Code slash commands.}
\label{tab:claude-code-slash}
\begin{tabular}{ll}
\toprule
\textbf{Command} & \textbf{Effect} \\
\midrule
\texttt{/help}            & List all available commands and skills \\
\texttt{/model [name]}    & Switch the active model \\
\texttt{/compact}         & Compress conversation history to free context \\
\texttt{/clear}           & Start a fresh session (context wiped) \\
\texttt{/config}          & Open interactive settings editor \\
\texttt{/permissions}     & Add or remove tool allowlist entries \\
\texttt{/memory}          & View and edit auto-memory files \\
\texttt{/status}          & Show current session token usage \\
\texttt{/fast}            & Toggle Fast mode (Opus with faster output) \\
\texttt{/terminal-setup}  & Write IDE keybindings for multi-line input \\
\texttt{/init}            & Generate a \texttt{CLAUDE.md} for the current repo \\
\texttt{/workflows}       & Show live progress of running Workflow tasks \\
\bottomrule
\end{tabular}
\end{table}

Any skill installed in \texttt{.claude/skills/} is also accessible as a slash
command (e.g., \texttt{/draft-chapter}, \texttt{/score-content}).

\subsection{Plan Mode}
\label{app:claude-code:plan-mode}

Plan mode is a read-only exploration mode in which Claude Code analyzes the
codebase, reasons about architecture, and produces a structured plan---but
cannot execute any writes or shell commands.  It is the safest way to understand
what a large, ambiguous task will require before granting execution permissions.

Activate Plan mode by pressing \textsc{Shift+Tab} once from the default Edit
mode; subsequent presses cycle through Auto-Accept mode and back.  The three
modes are:

\begin{enumerate}
  \item \textbf{Edit mode} (default) --- Claude proposes each file change
        individually; the user approves or rejects before writing.
  \item \textbf{Plan mode} --- read-only; Claude explores and writes a plan
        document; no files are modified.
  \item \textbf{Auto-Accept mode} --- Claude executes approved tool categories
        without per-call confirmation; see Section~\ref{app:claude-code:permissions}
        for how to configure which tools are auto-approved.
\end{enumerate}

\begin{deepdive}
In Plan mode, Claude Code internally activates the \texttt{EnterPlanMode} and
\texttt{ExitPlanMode} tools.  These tools are deferred-schema tools: their
parameter schemas are not loaded until explicitly requested, which is why Plan
mode imposes essentially zero overhead on sessions that never use it.  This
pattern---load schemas on demand---is used broadly in Claude Code to keep the
system prompt size bounded even when hundreds of extension tools are registered.
\end{deepdive}

\subsection{Keyboard Shortcuts}
\label{app:claude-code:shortcuts}

The most productive keyboard shortcuts differ slightly between macOS and
Windows/Linux.  Table~\ref{tab:claude-code-keys} gives the canonical bindings.

\begin{table}[ht]
\centering
\caption{Claude Code keyboard shortcuts (macOS / Windows--Linux).}
\label{tab:claude-code-keys}
\begin{tabular}{lll}
\toprule
\textbf{Action} & \textbf{macOS} & \textbf{Windows/Linux} \\
\midrule
Submit prompt (single line)    & \textsc{Return}    & \textsc{Return} \\
Newline in prompt              & \textsc{Shift+Return} & \textsc{Shift+Return} \\
Cycle Edit → Plan → Auto       & \textsc{Shift+Tab} & \textsc{Shift+Tab} \\
Open plan file in editor       & \textsc{Cmd+G}     & \textsc{Ctrl+G} \\
Undo last AI change            & \textsc{Esc, Esc}  & \textsc{Esc, Esc} \\
Toggle extended thinking       & \textsc{Option+T}  & \textsc{Alt+T} \\
Toggle editor/Claude focus     & \textsc{Cmd+Esc}   & \textsc{Ctrl+Esc} \\
Interrupt running agent        & \textsc{Ctrl+C}    & \textsc{Ctrl+C} \\
\bottomrule
\end{tabular}
\end{table}

The double-\textsc{Esc} rewind is particularly useful during iterative
drafting: it undoes all file changes from the last assistant turn and removes
that turn from the conversation history, allowing a clean retry with a different
prompt.

\subsection{Context Compaction and Conversation Management}
\label{app:claude-code:context}

Claude Code's context window grows with every tool call and assistant response.
When the context approaches its limit, Claude Code automatically runs
\texttt{/compact}: it summarises the conversation history into a compressed
representation that preserves task state while freeing tokens for continued
work.

Key points for practitioners:

\begin{itemize}
  \item Running \texttt{/compact} \emph{manually} before starting a new phase
        of work gives more predictable results than waiting for the automatic
        trigger.
  \item The Anthropic prompt cache has a 5-minute TTL.  Sleeping longer than
        300 seconds between turns (e.g., during a manual review step) incurs a
        cache miss; in the Workflow tool the scheduler accounts for this when
        choosing polling intervals.
  \item For very long autonomous runs, consider splitting the task into
        independent Workflow phases; each phase starts with a fresh agent and a
        focused prompt rather than a compacted summary of a long history.
\end{itemize}

% ---------------------------------------------------------------
\section{Configuration and Memory}
\label{app:claude-code:config}

\subsection{The CLAUDE.md File Hierarchy}
\label{app:claude-code:claude-md}

\texttt{CLAUDE.md} is the primary instruction file: Claude Code reads it
automatically at the start of every session and treats its content as
highest-priority guidance, overriding default system-prompt behaviour.  The
file system supports a layered hierarchy:

\begin{enumerate}
  \item \textbf{Global} (\texttt{\textasciitilde/.claude/CLAUDE.md}) ---
        applies to all projects on the machine; a good place for personal
        preferences and cross-project conventions.
  \item \textbf{Project} (\texttt{.claude/CLAUDE.md} at the repo root) ---
        applies to the current project; checked into version control so all
        collaborators (human and AI) see the same rules.
  \item \textbf{Sub-directory} (any \texttt{CLAUDE.md} in a nested folder) ---
        automatically included when Claude Code works within that directory;
        useful for monorepos with per-package conventions.
\end{enumerate}

All three levels are loaded and merged, with the more specific file taking
precedence on any conflict.  The \texttt{/init} command generates a starter
\texttt{CLAUDE.md} by reading the current directory structure, recent git
history, and detected languages/frameworks.

\begin{deepdive}
\texttt{CLAUDE.md} is not limited to free-form prose.  It can contain structured
YAML front matter (as in this book's \texttt{TOPIC.md}), conditional blocks, and
even embedded skill references.  Claude Code does not execute front matter
directly; rather, skills and agents read it programmatically using the
\texttt{Read} tool to extract configuration.  This indirection---human-readable
markdown that is also machine-parseable---is what allows a single configuration
file to govern both AI sessions and the shell hooks that validate them.
\end{deepdive}

\subsection{Settings, Permissions, and Environment Variables}
\label{app:claude-code:settings}

Runtime behaviour is controlled by JSON settings files at three scopes:

\begin{minted}{json}
// ~/.claude/settings.json       (global — lowest priority)
// .claude/settings.json         (project — medium priority)
// .claude/settings.local.json   (session — highest priority; gitignored)
\end{minted}

The most important keys are:

\begin{minted}{json}
{
  "model": "claude-sonnet-4-6",
  "permissions": {
    "allow": [
      "Bash(git log:*)",
      "Bash(npm test:*)",
      "Edit",
      "Read"
    ],
    "deny": [
      "Bash(rm -rf:*)",
      "Bash(curl:*)"
    ]
  },
  "env": {
    "ANTHROPIC_API_KEY": "${ANTHROPIC_API_KEY}",
    "PYTHONPATH": "${workspaceRoot}/src"
  },
  "hooks": {
    "PostToolUse": [
      { "matcher": "Edit", "command": "bash .claude/hooks/lint.sh" }
    ]
  }
}
\end{minted}

The \texttt{permissions.allow} list uses a \texttt{Tool(pattern:*)} syntax
where the pattern is matched against the tool's first argument.  The
\texttt{permissions.deny} list takes precedence over \texttt{allow}: if a
command matches both, it is denied.  This asymmetric design ensures that safety
constraints cannot be accidentally overridden by a broad allowlist entry.

The \texttt{/permissions} slash command provides an interactive editor for the
allowlist that does not require manually editing JSON.

\textbf{Environment variables.}  Variables declared in \texttt{settings.json}
under \texttt{env} are injected into every Bash tool call.  Use
\texttt{\$\{VAR\}} syntax to forward a variable from the parent shell.  Never
hardcode API keys in settings files that are committed to version control; use
\texttt{settings.local.json} (gitignored by default) or OS-level secret stores
instead.

\subsection{Auto-Memory System}
\label{app:claude-code:memory}

Claude Code maintains a persistent memory store at
\texttt{\textasciitilde/.claude/projects/\textless{}encoded-path\textgreater/memory/}.
Each memory is a markdown file with YAML front matter:

\begin{minted}{markdown}
---
name: feedback-no-paragraphs
description: Never use \paragraph{} as inline labels; prose must flow
metadata:
  type: feedback
---

Rule: Do not use `\paragraph{}` or `\textbf{}` as inline paragraph labels.
**Why:** The author finds heading-dense prose easier to scan in source but
harder to read in the compiled PDF.
**How to apply:** In all `.tex` content files, let sections flow as continuous
prose without sub-headings smaller than `\subsubsection`.
\end{minted}

The memory system has four types: \textbf{user} (the human's role and
preferences), \textbf{feedback} (corrections and confirmations), \textbf{project}
(goals, deadlines, decisions), and \textbf{reference} (pointers to external
systems).  A \texttt{MEMORY.md} index file is always loaded into context; the
full memory files are read on demand when flagged as relevant by the index.

Use the \texttt{/memory} command to inspect and edit the memory store from
within a Claude Code session.

% ---------------------------------------------------------------
\section{Built-in Tools}
\label{app:claude-code:tools}

Claude Code exposes a curated set of tools to the language model.  Each tool
call requires user approval unless the tool (or a matching pattern) is listed in
\texttt{permissions.allow}.

\subsection{File System Tools}
\label{app:claude-code:tools-fs}

\begin{description}
  \item[\texttt{Read}] Reads a file at an absolute path.  Supports line-range
        offsets (\texttt{offset}, \texttt{limit}) for large files; for PDFs
        accepts a \texttt{pages} range.  Always read a file before editing it:
        the \texttt{Edit} tool refuses to operate on files it has not seen.

  \item[\texttt{Write}] Overwrites a file with new content.  Use only for new
        files or complete rewrites; prefer \texttt{Edit} for partial changes.

  \item[\texttt{Edit}] Applies an exact string replacement (\texttt{old\_string}
        → \texttt{new\_string}) in a file.  The \texttt{replace\_all} flag
        performs the substitution everywhere in the file, useful for renaming a
        symbol throughout a module.  The replacement fails if \texttt{old\_string}
        is not unique, preventing accidental edits.

  \item[\texttt{Glob}] Fast file-pattern matching.  Returns paths sorted by
        modification time.  Use instead of \texttt{find} for any file-discovery
        task.

  \item[\texttt{Grep}] Content search backed by \texttt{ripgrep}.  Supports
        full regex syntax, file-type filtering (\texttt{type} parameter),
        multi-line matching, and context lines (\texttt{-C}).  Output modes:
        \texttt{files\_with\_matches}, \texttt{content}, or \texttt{count}.
\end{description}

\subsection{The Bash Tool}
\label{app:claude-code:tools-bash}

The \texttt{Bash} tool executes arbitrary shell commands with a configurable
timeout (up to 600\,s).  Working directory persists across calls within a
session; environment state (variables, functions) does not.

Commands can be run in the background by setting \texttt{run\_in\_background:
true}, in which case Claude Code is notified when the process completes rather
than blocking.  This is essential for long compilation jobs or test suites.

\begin{minted}{bash}
# Example: compile the book in the background
cd book && pdflatex -interaction=nonstopmode main.tex
\end{minted}

\textbf{Security note.}  The Bash tool is the highest-risk tool because it can
execute arbitrary code.  Use the \texttt{permissions.deny} list to block
patterns that should never run in your environment (e.g., \texttt{curl} calls
to external hosts, \texttt{rm -rf} without path specifics).

\subsection{Orchestration Tools}
\label{app:claude-code:tools-orchestration}

Three tools control multi-agent execution:

\begin{description}
  \item[\texttt{Agent}] Spawns a child Claude Code session with its own context
        window, prompt, and tool set.  Returns the agent's final text output (or
        a structured JSON object if a \texttt{schema} is provided).  Useful for
        parallelising independent tasks or protecting the main context window
        from excessive search results.  The optional
        \texttt{isolation: "worktree"} flag runs the agent in a fresh git
        worktree, preventing file conflicts when multiple agents write in
        parallel.

  \item[\texttt{Skill}] Invokes a named skill (slash command) from
        \texttt{.claude/skills/} and returns its output.  Equivalent to typing
        \texttt{/skill-name} at the prompt, but callable programmatically from
        within another skill or Workflow script.

  \item[\texttt{Workflow}] Executes a JavaScript orchestration script that fans
        out work across many agents deterministically.  See
        Section~\ref{app:claude-code:workflow} for the full API.
\end{description}

\subsection{Extended Thinking}
\label{app:claude-code:thinking}

Claude Opus models support \emph{extended thinking}: an internal chain-of-thought
scratchpad in which the model reasons before producing output.  Claude Code
exposes this via the \texttt{Alt+T} / \texttt{Option+T} keyboard shortcut or
the \texttt{/fast} toggle (which disables thinking in exchange for lower
latency).

Extended thinking is most valuable for:
\begin{itemize}
  \item architectural decisions that span many files,
  \item mathematical derivations that require multi-step symbolic reasoning, and
  \item debugging sessions where the root cause is non-obvious.
\end{itemize}

When extended thinking is active, the model's reasoning tokens are \emph{not}
billed at the same rate as output tokens; consult the Anthropic pricing page for
current rates.

% ---------------------------------------------------------------
\section{Extension Mechanisms}
\label{app:claude-code:extensions}

\subsection{Skills}
\label{app:claude-code:skills}

A \emph{skill} is a markdown file at
\texttt{.claude/skills/\textless{}name\textgreater{}/SKILL.md} that describes a
multi-step workflow.  When the user types \texttt{/\textless{}name\textgreater{}}
at the prompt, Claude Code loads the file and follows its instructions exactly.
Skills are:

\begin{itemize}
  \item \textbf{Portable} --- the same \texttt{SKILL.md} file works in Claude
        Code, Codex CLI, Gemini CLI, and Cursor without modification.
  \item \textbf{Version-controlled} --- stored in the repo alongside the code
        they govern; changes are visible in \texttt{git log}.
  \item \textbf{Composable} --- a skill can invoke other skills via the
        \texttt{Skill} tool, enabling pipelines like
        \texttt{/draft-chapter} → \texttt{/score-content} →
        \texttt{/refine-until-threshold}.
\end{itemize}

This book defines skills for every major authoring task; see the
\texttt{.claude/skills/} directory for the full catalogue, or run
\texttt{/help} to list available commands.

\subsection{Hooks}
\label{app:claude-code:hooks}

Hooks are shell scripts registered in \texttt{settings.json} that execute
automatically in response to tool-use events.  The available hook events are:

\begin{description}
  \item[\texttt{PreToolUse}] Fires before a tool call; can cancel the call by
        exiting non-zero.
  \item[\texttt{PostToolUse}] Fires after a tool call completes; its output is
        shown to the model as feedback.
  \item[\texttt{UserPromptSubmit}] Fires when the user submits a prompt; can
        prepend additional context.
  \item[\texttt{Stop}] Fires when Claude Code finishes a turn; useful for
        notifications or cleanup.
\end{description}

\begin{minted}{json}
"hooks": {
  "PostToolUse": [
    {
      "matcher": "Edit",
      "command": "bash .claude/hooks/lint-on-save.sh \"$CLAUDE_TOOL_INPUT_FILE\""
    }
  ],
  "Stop": [
    { "command": "bash .claude/hooks/notify.sh" }
  ]
}
\end{minted}

Hook scripts receive tool input and output as environment variables and can
write to \texttt{stdout} to inject messages into the conversation.  A common
pattern in this book's project is a \texttt{PostToolUse} hook on \texttt{Edit}
that re-runs the \LaTeX{} linter on any modified \texttt{.tex} file,
surfacing syntax errors before the next agent step.

\subsection{MCP Server Integration}
\label{app:claude-code:mcp}

The Model Context Protocol (MCP) is an open standard (released November 2024)
that allows Claude Code to connect to external processes that expose tools and
data sources.  An MCP server is a long-running process that Claude Code discovers
via \texttt{settings.json}:

\begin{minted}{json}
{
  "mcpServers": {
    "context7": {
      "command": "npx",
      "args": ["-y", "@upstash/context7-mcp@latest"]
    },
    "sqlite": {
      "command": "uvx",
      "args": ["mcp-server-sqlite", "--db-path", "data/research.db"]
    }
  }
}
\end{minted}

MCP servers expose their tools as callable functions inside Claude Code's tool
list.  The finance-relevant servers include database connectors (SQLite,
PostgreSQL), browser automation (for scraping financial filings), internal API
bridges (Bloomberg terminal, FactSet), and document stores (Google Drive, email).

Tools from MCP servers are deferred: their parameter schemas are not loaded into
context until the user invokes a \texttt{ToolSearch} query, keeping the system
prompt compact when many servers are registered.

\subsection{Plugins}
\label{app:claude-code:plugins}

A \emph{plugin} (introduced October 2025) is a packaged bundle that can include
any combination of skills, MCP server configurations, sub-agent definitions, and
hooks.  Plugins are installed with:

\begin{minted}{bash}
claude plugin install @anthropic/financial-data-plugin
\end{minted}

After installation, all of the plugin's components are available without manual
configuration.  The community plugin registry hosts hundreds of domain-specific
plugins; for regulated environments, private registries can be pointed to via
the \texttt{pluginRegistry} setting.

% ---------------------------------------------------------------
\section{Agents and Multi-Agent Systems}
\label{app:claude-code:agents}

\begin{context}
The move from a single LLM call to a \emph{system} of coordinated LLM calls is
arguably the most consequential architectural shift in applied AI since the
Transformer itself.  A single model, no matter how capable, operates within a
fixed context window; a system of models can divide a large task into pieces,
process them in parallel, check each other's work, and synthesise results that
no single context window could hold.  Claude Code's multi-agent architecture is
designed around exactly this observation.
\end{context}

\subsection{Subagents and the Agent Tool}
\label{app:claude-code:subagents}

The \texttt{Agent} tool spawns a child Claude Code session.  Each child has:

\begin{itemize}
  \item its own independent context window,
  \item a prompt written by the parent agent,
  \item a configurable tool set (defaulting to all tools), and
  \item an optional \texttt{schema} that forces the child to return a validated
        JSON object rather than free text.
\end{itemize}

Children run concurrently when multiple \texttt{Agent} calls appear in the same
response turn.  The parent collects results and synthesises them.  The most
common pattern is a two-level hierarchy: a \emph{planner} agent that breaks the
task into pieces, followed by a fleet of \emph{worker} agents that execute each
piece in parallel.

\begin{minted}{json}
// Simplified Agent call (as JSON schema)
{
  "description": "Classify whether chapter section covers tokenization",
  "prompt": "Read book/chapters/02-llm-foundations/chapter.tex and return ...",
  "schema": {
    "type": "object",
    "properties": {
      "covers_tokenization": { "type": "boolean" },
      "evidence": { "type": "string" }
    }
  }
}
\end{minted}

\subsection{Custom Subagent Types}
\label{app:claude-code:custom-agents}

Named sub-agent types are defined as markdown files in
\texttt{.claude/agents/}.  Each file specifies a persona, a set of inputs, a
task description, and constraints.  When an \texttt{Agent} call includes
\texttt{agentType: "math-checker"}, Claude Code reads
\texttt{.claude/agents/math-checker.md} and prepends it to the child's system
prompt.

This book defines the following agent types:

\begin{itemize}
  \item \texttt{book-writer} --- drafts \LaTeX{} content following the book's
        structural conventions.
  \item \texttt{math-checker} --- verifies derivations and proofs; returns
        \texttt{PASS} or \texttt{FAIL} with a list of issues.
  \item \texttt{editor} --- improves clarity and flow without changing
        mathematical content.
  \item \texttt{scorer} --- applies the 5-dimension quality rubric and writes
        a JSON score report.
  \item \texttt{outline-curator} --- audits and refines section outlines.
  \item \texttt{literature-reviewer} --- suggests references with BibTeX entries.
  \item \texttt{proofreader} --- checks grammar, spelling, and style consistency.
\end{itemize}

\subsection{The Workflow Tool}
\label{app:claude-code:workflow}

The \texttt{Workflow} tool runs a JavaScript orchestration script that fans
work out across many agents deterministically.  Unlike the \texttt{Agent} tool
(which is invoked by the model and subject to the model's discretion), a
Workflow script is controlled by explicit JavaScript control flow---loops,
conditionals, and fan-out---making the execution pattern reproducible and
auditable.

A minimal Workflow script:

\begin{minted}{javascript}
export const meta = {
  name: 'score-all-chapters',
  description: 'Run quality scorer on every chapter in parallel',
  phases: [{ title: 'Score' }],
}

const chapters = ['01-intro', '02-llm-foundations', '03-llm-training-finetuning']

const SCORE_SCHEMA = {
  type: 'object',
  properties: {
    clarity: { type: 'number' },
    rigor: { type: 'number' },
    completeness: { type: 'number' },
    pedagogy: { type: 'number' },
    style: { type: 'number' },
  },
}

const results = await parallel(chapters.map(ch => () =>
  agent(
    `Score book/chapters/${ch}/chapter.tex on 5 dimensions (0-10 each).`,
    { label: `score:${ch}`, phase: 'Score', schema: SCORE_SCHEMA }
  )
))

return results
\end{minted}

Key Workflow primitives:

\begin{description}
  \item[\texttt{agent(prompt, opts)}] Spawns one sub-agent; returns the
        agent's text output or, when \texttt{opts.schema} is provided, a
        validated JSON object.
  \item[\texttt{parallel(thunks)}] Runs all thunk functions concurrently;
        waits for all to complete (barrier); returns an array of results.
  \item[\texttt{pipeline(items, ...stages)}] Runs each item through all
        stages without inter-stage barriers.  The default choice: use
        \texttt{parallel} only when all prior-stage results are genuinely
        needed before the next stage can start.
  \item[\texttt{phase(title)}] Opens a named progress group; subsequent
        \texttt{agent} calls are grouped under this title in the live
        \texttt{/workflows} view.
  \item[\texttt{log(message)}] Emits a progress message visible in the UI.
\end{description}

Workflows support \emph{resume}: if a run is interrupted, re-launching with
the same \texttt{scriptPath} and \texttt{resumeFromRunId} replays cached
agent results and continues from the first changed or new call.

\subsection{Agent Teams}
\label{app:claude-code:agent-teams}

Agent Teams (released February 2026) extend the parent/child subagent model
to a \emph{peer} coordination model.  A team consists of:

\begin{description}
  \item[Team Lead] The main Claude Code session that initialises the team,
        posts the initial task list, and synthesises final results.
  \item[Teammates] Independent Claude Code processes, each with their own
        context window.  Teammates can read from and write to a shared task
        list, and can send direct messages to each other via a mailbox system.
\end{description}

Agent Teams require Claude Opus 4.8 for all participants and are disabled by
default.  Enable them by setting the environment variable
\texttt{CLAUDE\_CODE\_EXPERIMENTAL\_AGENT\_TEAMS=1}.  They are most productive
for tasks that benefit from genuine collaboration: parallel code review, complex
debugging where teammates test competing hypotheses, or large feature
implementations where different teammates own different layers (frontend,
backend, database, tests).

Token consumption is roughly 7$\times$ a single-session Plan-mode run for
comparable tasks; reserve Agent Teams for problems that would otherwise require
repeated manual context-switching.

% ---------------------------------------------------------------
\section{Automation and Remote Execution}
\label{app:claude-code:automation}

\subsection{Scheduled Routines}
\label{app:claude-code:routines}

As of March 2026, Claude Code supports \emph{Routines}: cron-style tasks that
execute on Anthropic-managed infrastructure, decoupled from the developer's
local machine.  Create a routine with the \texttt{/schedule} command:

\begin{minted}{text}
/schedule "every Monday 08:00 UTC: /full-review all chapters, commit if PASS"
\end{minted}

Routines are stored in \texttt{.claude/routines/} and have the same access to
project skills and MCP servers as an interactive session.  Common use cases for
research workflows:

\begin{itemize}
  \item Nightly bibliography audits that flag broken DOI links.
  \item Weekly quality re-scores after upstream dependency updates.
  \item Post-merge hooks that re-sync lecture slides with their companion
        chapters whenever a chapter \texttt{.tex} file changes.
\end{itemize}

\subsection{External Integrations}
\label{app:claude-code:integrations}

Claude Code integrates with Slack as a first-party surface.  After authorising
the integration, mentioning \texttt{@Claude} in a Slack channel or thread
triggers a Claude Code session on Anthropic infrastructure; the session has
access to the linked repository and can commit results back to GitHub.

For CI/CD integration, Claude Code can be invoked non-interactively:

\begin{minted}{bash}
# In a GitHub Actions workflow step:
claude --non-interactive \
       --model claude-sonnet-4-6 \
       "Run /score-content on all modified .tex files in this PR. \
        Fail if any dimension is below 8."
\end{minted}

The \texttt{--non-interactive} flag disables the permission-approval UI; the
allowlist in \texttt{.claude/settings.json} governs what the agent is permitted
to do.

% ---------------------------------------------------------------
\section{Security, Privacy, and Cost}
\label{app:claude-code:security}

\subsection{The Permissions Model}
\label{app:claude-code:permissions}

Claude Code's permission system follows a \emph{default-cautious} principle:
every tool call that is not explicitly allowed prompts the user for approval.
The layered allowlist/denylist design described in
Section~\ref{app:claude-code:settings} provides fine-grained control.

For shared codebases, commit a minimal \texttt{.claude/settings.json} that
allows only the operations required for automated tasks (read-only tools plus
specific git operations) and requires human approval for anything that touches
production data or external network endpoints.

The \texttt{permissions.deny} key provides a hard veto.  Any pattern listed
there is blocked regardless of what the allowlist says or what the model
requests.  Finance teams should at minimum deny:

\begin{minted}{json}
"deny": [
  "Bash(curl:*)",
  "Bash(wget:*)",
  "Bash(pip install:*)",
  "Bash(npm install:*)",
  "Bash(rm -rf:*)"
]
\end{minted}

This prevents the model from inadvertently (or through a prompt-injection
attack) exfiltrating data or installing arbitrary packages.

\subsection{Sensitive Data and Secrets Management}
\label{app:claude-code:secrets}

Claude Code reads the project directory, which may contain data files, model
weights, or configuration with credentials.  Best practices:

\begin{enumerate}
  \item \textbf{Gitignore secrets} --- add \texttt{.env}, \texttt{*.pem},
        and \texttt{secrets.json} to \texttt{.gitignore} before running
        \texttt{/init}, which reads git history.
  \item \textbf{Deny read access} ---
        \texttt{"deny": ["Read(.env*)", "Read(*credentials*)"]} prevents
        the model from inadvertently including secret values in its context.
  \item \textbf{Use \texttt{settings.local.json}} --- environment variables
        injected via this file are not committed to the repository; use it
        for per-developer API keys.
  \item \textbf{Redact before sharing transcripts} --- the \texttt{/export}
        command saves the conversation; review the export before sharing, as
        it may capture tool outputs that include file contents.
\end{enumerate}

For financial data subject to NDA or data sharing agreements, consider running
Claude Code with a local model endpoint (via the \texttt{baseURL} setting)
rather than the Anthropic API, so data never leaves your perimeter.

\subsection{Token Consumption and Cost Control}
\label{app:claude-code:cost}

Claude Code sessions can consume substantial tokens, particularly when extended
thinking is active or when Workflow scripts run large agent fleets.
Table~\ref{tab:token-costs} gives approximate consumption for the tasks most
relevant to this book.

\begin{table}[ht]
\centering
\caption{Approximate token consumption for common Claude Code tasks (Sonnet 4.6).}
\label{tab:token-costs}
\begin{tabular}{lrr}
\toprule
\textbf{Task} & \textbf{Input tokens} & \textbf{Output tokens} \\
\midrule
\texttt{/score-content} on one chapter & 15{,}000--30{,}000 & 2{,}000--4{,}000 \\
\texttt{/draft-chapter} full pipeline & 50{,}000--120{,}000 & 10{,}000--25{,}000 \\
\texttt{/full-review} (all agents)     & 80{,}000--200{,}000 & 15{,}000--40{,}000 \\
Agent Teams run (7 sessions)           & 350{,}000+          & 70{,}000+ \\
\bottomrule
\end{tabular}
\end{table}

Cost-reduction strategies:

\begin{itemize}
  \item Use Sonnet 4.6 for bulk operations (scoring, proofreading) and reserve
        Opus 4.8 for tasks requiring deep reasoning.
  \item Exploit prompt caching: CLAUDE.md and long system prompts are cached by
        Anthropic with a 5-minute TTL; rapid iteration within a session benefits
        automatically.
  \item Set the \texttt{budget} directive in Workflow scripts
        (\texttt{+500k tokens}) to cap a run and prevent runaway agent loops.
  \item Run \texttt{/compact} at natural task boundaries rather than letting
        the context grow until automatic compaction triggers.
\end{itemize}

% ---------------------------------------------------------------
\section{Claude Code in Research and Finance Workflows}
\label{app:claude-code:finance}

\subsection{Use Cases for Quants and Researchers}
\label{app:claude-code:use-cases}

The following use cases have proven especially productive for the audience of
this book:

\begin{description}
  \item[Literature synthesis] A Workflow script that fans out across a list of
        SSRN or arXiv paper IDs, has one sub-agent per paper extract
        methodology and key results as structured JSON, then has a synthesis
        agent produce a comparative table---in minutes rather than days.

  \item[Code archaeology] Dropping a legacy pricing library into the project
        directory and asking Claude Code to generate a call graph, identify
        undocumented assumptions, and produce a test suite covering edge cases.

  \item[Report drafting] Using \texttt{/draft-chapter} and
        \texttt{/score-content} iteratively to produce investment memos or
        regulatory reports that meet a house style and a readability threshold.

  \item[Data pipeline debugging] Providing a failing notebook and asking Claude
        Code to reproduce the error, hypothesise causes, write a minimal
        reproducer, and propose a fix---all within one non-interactive
        \texttt{claude} invocation.

  \item[Nightly model monitoring] A scheduled routine that runs a battery of
        sanity checks on a live NLP model (distribution shift, calibration, new
        OOV tokens), commits a JSON report, and posts a summary to Slack if any
        check fails.
\end{description}

\subsection{This Book as a Working Example}
\label{app:claude-code:book}

This book was produced using Claude Code as the primary authoring environment.
The \texttt{.claude/} directory at the repository root encodes the entire
production workflow:

\begin{itemize}
  \item \texttt{.claude/CLAUDE.md} --- master instructions: conventions,
        quality gates, commit message format.
  \item \texttt{.claude/skills/} --- one skill per workflow step
        (scaffold, draft, score, refine, release).
  \item \texttt{.claude/agents/} --- specialised personas (book-writer,
        math-checker, scorer, etc.).
  \item \texttt{.claude/hooks/} --- shell scripts that lint \LaTeX{} on save
        and re-run quality scores on pre-commit.
  \item \texttt{docs/quality/} --- JSON score reports for every chapter,
        committed alongside the content.
  \item \texttt{TOPIC.md} --- YAML configuration for subject, audience,
        quality threshold, and auto-commit behaviour.
\end{itemize}

Replicating this structure for a new research project requires running
\texttt{/interview-me} to configure \texttt{TOPIC.md}, then
\texttt{/new-topic} for each chapter.  The entire scaffold is version-controlled
and portable: cloning the repository gives any collaborator (human or AI) the
same working environment.

\subsection{Recommended Workflow Patterns}
\label{app:claude-code:patterns}

\begin{enumerate}
  \item \textbf{Plan before writing.}  Use Plan mode to review the proposed
        changes before committing to Auto-Accept.  For new chapters, run
        \texttt{/new-topic} in Plan mode first; approve the outline, then
        switch to Edit mode for the actual scaffold.

  \item \textbf{Score early, score often.}  Run \texttt{/score-content} after
        every significant draft, not just at the end.  The quality JSON
        provides a compass for the next iteration that is more reliable than
        subjective impression.

  \item \textbf{Separate concerns with agents.}  Do not ask a single session to
        draft content, check math, and improve style in the same turn.  Each
        task is well-suited to a different agent persona; the \texttt{Agent}
        tool makes this separation cheap.

  \item \textbf{Use the deny list proactively.}  In research environments with
        sensitive data, configure \texttt{permissions.deny} before the first
        session, not after a concern arises.

  \item \textbf{Compact at phase boundaries.}  Run \texttt{/compact} when
        transitioning from drafting to reviewing, or from reviewing to
        releasing.  This preserves the task state while eliminating draft
        history that is no longer actionable.

  \item \textbf{Commit checkpoints.}  The skills in this book commit after every
        major step.  Frequent small commits make it easy to bisect a regression
        introduced by an agent and revert to a known-good state without losing
        days of work.
\end{enumerate}
